\section{A measurement failure before offline optimization}
\label{sec:offline-scoring-abort}
The RTX3090 receiver passed the independent reviewed numerical audit, but the
subsequent offline training experiment did not complete. Frozen S+R+F scoring
closed 211 of the 1,000 new prompt groups before a rhythm measurement failed.
The completed prefix contained 844 candidates, of which 516 were admitted and
328 were rejected; eight prefix groups had no admitted candidates. These are
ordered-prefix statistics, not an unbiased complete-cache estimate or a
discrimination-accuracy benchmark.

The full-cache CPU quality check identified 24 groups whose four candidates
failed unchanged pre-analyzer quality gates. Combined with the known scored
rejections and assuming all remaining groups and all 50 older groups were
usable, the conservative signal-group upper bound was 1,025. Full merge and
eligibility preflight remained incomplete; neither feasibility nor insufficient
data for the fixed 1,000-update target was established.

The immediate exception was \texttt{ValueError: Not all downbeats are beats.}
from BeatThis's beat/downbeat export. The subprocess worker exited nonzero and
the strict bridge converted its failure to \texttt{RuntimeError}, bypassing the
reward bridge's existing measurement-invalid \texttt{ValueError} handling.
This terminated scoring before any new optimizer update. There was no new RL
checkpoint and no new fixed-test official/RL audio comparison. The precise
candidate and underlying beat-array discrepancy require controlled reproduction.

All 314 archived output members were reverified by SHA-256. The 42 original
input archives, totaling 132,352,450,560 bytes, remained on persistent storage;
the source audio/trajectory cache was not copied to the Mac. Protected closeout
released the instance and official reads confirmed stopped compute billing,
with a final bill of CNY9.86. However, releasing immediately on a terminal stage
error was too aggressive for this recoverable measurement failure; future runs
need a bounded diagnosis/recovery window before protected shutdown.

A separate opt-in error adapter was prepared after release. It recognizes only
the observed traceback, pins the original bridge/worker identities, rechecks
the original runtime status and produces no synthetic feature or valid score.
Twelve CPU contract tests passed, but no new GPU/audio validation or production
activation occurred. Full scoring, 1,000 additional actual updates and the
50-prompt held-out comparison remain future work rather than claimed results.

Source, raw-output and analysis references are given in
\texttt{REPORT\_ABORTED\_SCORING\_20261031.md}; the local artifact packet is
\texttt{music/rl\_results/offline\_grpo1000\_v2\_20261031\_2128889} and the
persistent packet is
\texttt{/mnt/ai-music-offline-grpo1000-v2-20261031-2128889}.
The recovery adapter is \texttt{src/music\_detector/rl/measurement\_rejections.py},
with contract tests in \texttt{tests/test\_rl\_measurement\_rejections.py}.
Raw terminal output is retained in the lossless \texttt{work.tar.gz.part000}
archive; derived statistics are \texttt{SCORING\_FAILURE\_ANALYSIS.json} and
\texttt{PREFIX\_GROUP\_SCORES.csv}. Protected-release and billing evidence is
referenced by \texttt{FINAL\_CLOSEOUT.json}. These distinguish executable code,
original outputs, derived analysis and resource-closeout evidence.
